\section*{Exponentially accurate randomized circuits}
The same gate coefficient holds at two-sided pointwise error
$2^{-eN}$ whenever
\[
 e>d+(D-2d)_+a/q.
\]
As in the theorem, assume $\kappa=d+(D-2d)_+a/q<b$. This gives,
for example, the strict $2.5226N$ bound at the first threshold in
the table for every fixed $e\ge .00564$. The same exponent controls uniform false positives and false negatives
on the common easy distribution.

To prove this claim, fix a gate coefficient strictly below the
theorem's bound. Choose constants $a<a'<a''$ sufficiently close to
$a$. Use $a''$ in the twig cutoff and write
$\kappa''=d+(D-2d)_+a''/q$. All strict margins, including
$b>\kappa''$ and $e>\kappa''$, persist. For any
deterministic circuit respecting the new gate bound, the two
concentration estimates in the proof imply that the maximum restored
dimension is at least $a'N$ with tilted probability $1-o(1)$.
This inference only uses the gate count. In particular it holds for
circuits with arbitrary erroneous labels. When the cutoff is not
reached the same algebra applies with $a'$ in place of $a$, since
otherwise the size lower bound with $a'$ would contradict our strict
gate margin.

Consider the generalized terminal path that records each affine
answer and each twig \emph{output}. Choose simplifications from these
answers alone. The likelihood ratio relative to uniform is the
product of the ordinary branch likelihood ratios and the factors
$4q/3$ or $4(1-q)$ at tilted twigs. Its logarithm has bounded
increments, with conditional means $d$ and $D$, respectively.
The pathwise bound on their sum is at most $\kappa''N+O(1)$.
Center the increments, pad to $N$ steps, and apply Azuma's inequality.
Except on a set of tilted mass $o(1)$, the logarithm is at most
$\kappa''N+o(N)$. Thus the set of generalized paths of restored
dimension at least $a'N$ has uniform mass at least
$2^{-(\kappa''+o(1))N}$.

There is an exact fractional representation of each such generalized
path by affine certificates. In a common twig let $H$ be its three
possible pairs. Any two distinct binary points constitute an affine
line. Each point of $H$ belongs to two of its three two-point lines, so
\[
 \mathbf 1_H=\tfrac12\sum_{\ell\subset H,\ |\ell|=2}\mathbf 1_\ell .
\]
The third line, in addition to those that fix one forcing literal,
fixes the parity of the two bits. For each common twig choose one
of these three line equations and retain both literal equations at
exceptional twigs. Tensoring the displayed identity over the $m$
common twigs expresses the path indicator as $2^{-m}$ times the sum
of the $3^m$ resulting certificate indicators. The coordinates used
in each twig are free parameters at that stage, both removed before
later substitutions, so their affine pullbacks give one independent
restored direction per common twig for every choice of lines. In
particular every certificate has the full restored dimension of its
generalized path. This identity permits overlap with its exact
pointwise multiplicity.

Use the dense affine law of the companion manuscript \cite{local}
with threshold $c_1N/n$ and rank slightly above $(1-a')N$.
Its mass in every constituent certificate of a large-dimension path
is at least $1-o(1)$ times the uniform mass there. The positive
weighted identity transfers the same inequality to the entire
generalized path. Let $u$ denote uniform false-positive mass
on inputs of complexity above $c_2N/n$, and let $v$ denote the
false-negative mass under this common easy law. Large-dimension
paths labeled one have uniform mass at most $u+2^{-(b-o(1))N}$
by the census. Those labeled zero have uniform mass at most
$(1+o(1))v$ by fractional domination and disjointness of the
generalized paths. Hence
\[
 u+2^{-(b-o(1))N}+(1+o(1))v
 \ \ge\ 2^{-(\kappa''+o(1))N}.
\]
All asymptotic error bounds are uniform in the fixed circuit.
Averaging over independent circuit coins now gives the claimed
two-sided confidence statement.

\paragraph{A linear dependence on the density deficit at bounded error.}
For fixed $0<c<1$ and threshold $\lfloor cN/n\rfloor$, a bounded-error
outer circuit requires at least $(4c-3/2-o(1))N$ gates. More generally
this conclusion uses the completeness density $c=c_1$ in the promise,
and any fixed soundness density $c_2<1$. Here the two ensembles give
asymptotically identical output distributions for every smaller gate
coefficient.

Indeed use ordinary fair branches and fair raw twig pairs, so the point
law is uniform and every twig is common with probability $3/4$.
The concentration estimates in the proof give, with probability
$1-o(1)$,
\[
 M_{\rm initial}\ge\tfrac72 T+\tfrac12 W-o(N),\qquad
 L\ge \tfrac34 W-o(N),\qquad L\ge N-T-\tfrac14 W-o(N).
\]
For any fixed $a'>1-c$, if $L\le a'N$, substitution yields
$M_{\rm initial}\ge (7/2-4a')N-o(N)$. Since $M_{\rm initial}\le G+N$,
a gate coefficient strictly below $5/2-4(1-c)$ therefore makes
$L>a'N$ typical for some $a'>1-c$. The fractional identity gives
relative domination on every such generalized path. Their uniform
mass is $1-o(1)$, so summing the smaller path probabilities gives
$1-o(1)$ as well. The resulting total variation bound holds for
each fixed circuit and after averaging its coins. Pointwise error
bounded below $1/2$ by a fixed constant would require a positive
distinguishing advantage, since the uniform table is a required no
input with probability $1-o(1)$.
